\subsection{Diagonalizable Subalgebras of Simple Algebras}

Let D be a K-algebra that is a field, and let V be a finite-dimensional right vector space over D. Let L be a K-subalgebra of \(\operatorname{End}_{\mathrm{D}}(\mathrm{V})\) that is a diagonalizable K-algebra (V, §6, No.~3, p.~28). By definition, L has finite degree over K, and there exists a basis \((\varepsilon_{i})_{i\in\mathrm{I}}\) of L over K with the following properties:

\[
\varepsilon_ {i} ^ {2} = \varepsilon_ {i}, \varepsilon_ {i} \varepsilon_ {j} = 0 \text { if } i \neq j, \quad \sum_ {i \in I} \varepsilon_ {i} = 1.
\]

Set \(V_{i} = \varepsilon_{i}(V)\) for every i in I; then \((\mathrm{V}_{i})_{i \in \mathrm{I}}\) is a family of nonzero linear subspaces of V with direct sum V (II, §1, No.~8, p.~209, Proposition 12). Let u be an endomorphism of V; then u belongs to L if and and only if for every \(i \in I\) , there exists an element \(\lambda_{i}\) of K such that \(u(x) = \lambda_{i}x\) for every \(x \in V_{i}\) .

Conversely, suppose that V is the direct sum of a family \((\mathrm{V}_{i})_{i\in\mathrm{I}}\) of nonzero linear subspaces. For any element \(\boldsymbol{\lambda}=(\lambda_{i})_{i\in\mathrm{I}}\) of \(K^{I}\) , denote by \(u_{\lambda}\) the endomorphism of the D-vector space V such that \(u_{\boldsymbol{\lambda}}(x)=\lambda_{i}x\) for \(x\in V_{i}\) . The set L of endomorphisms \(u_{\lambda}\) for \(\lambda\in K^{I}\) is a diagonalizable subalgebra of \(\operatorname{End}_{\mathrm{D}}(\mathrm{V})\) having the family \((\varepsilon_{i})_{i\in\mathrm{I}}\) as basis, where \(\varepsilon_{i}\) is the projector with image \(V_{i}\) and kernel \(\sum_{j\neq i}V_{j}\) . We say that L is the diagonalizable subalgebra of \(\operatorname{End}_{\mathrm{D}}(\mathrm{V})\) associated with the direct sum decomposition \(V=\oplus_{i\in\mathrm{I}}V_{i}\) . We have \([L:K]=\operatorname{Card}(I)\leqslant\dim_{\mathrm{D}}(V)\) .

PROPOSITION 6. --- Let L be the diagonalizable subalgebra of \(\operatorname{End}_{\mathrm{D}}(\mathrm{V})\) associated with a direct sum decomposition \(V = \oplus_{i \in I} V_i\) .

\begin{enumerate}
\def\labelenumi{\alph{enumi})}
\item
  The algebra L is maximal among the diagonalizable subalgebras of the K-algebra \(\operatorname{End}_{\mathrm{D}}(\mathrm{V})\) if and only if every \(V_{i}\) has dimension 1 over D.
\item
  The algebra \(\mathrm{L}\) is a maximal commutative subalgebra of \(\operatorname{End}_{\mathrm{D}}(\mathrm{V})\) if and only if we have \(\mathrm{D} = \mathrm{K}\) and every \(\mathrm{V}_i\) has dimension 1 over \(\mathrm{K}\) .
\end{enumerate}

If every vector space \(V_{i}\) has dimension 1 over D, then we have

\[
[ \mathrm{L}: \mathrm{K} ] = \operatorname{Card} (\mathrm{I}) = \dim_ {\mathrm{D}} (\mathrm{V}),
\]

so L is maximal among the diagonalizable subalgebras of \(\operatorname{End}_{\mathrm{D}}(\mathrm{V})\) . In the opposite case, there exists an index \(j \in I\) such that \(\dim_{\mathrm{D}}(\mathrm{V}_{j}) \geqslant 2\) . Choose two nonzero linear subspaces \(V_{j}^{\prime}\) and \(V_{j}^{\prime\prime}\) of \(V_{j}\) with direct sum \(V_{j}\) . The diagonalizable subalgebra of \(\operatorname{End}_{\mathrm{D}}(\mathrm{V})\) associated with the direct sum decomposition \(\mathrm{V} = (\oplus_{i \in \mathbf{I} - \{j\}} \mathrm{V}_{i}) \oplus \mathrm{V}_{j}^{\prime} \oplus \mathrm{V}_{j}^{\prime\prime}\) contains L and is not equal to L; assertion a) follows.

The commutant \(L'\) of L in \(\operatorname{End}_{\mathrm{D}}(\mathrm{V})\) consists of the endomorphisms of the form \((x_{i}) \mapsto (u_{i}(x_{i}))\) , with \((u_{i}) \in \prod_{i \in \mathrm{I}} \operatorname{End}_{\mathrm{D}}(\mathrm{V}_{i})\) . The algebra L is a maximal commutative subalgebra of \(\operatorname{End}_{\mathrm{D}}(\mathrm{V})\) if and only if we have \(L = L'\) (VIII, p.~261, Lemma 3, a)). This relation is therefore equivalent to `` \(\operatorname{End}_{\mathrm{D}}(\mathrm{V}_{i}) = \mathrm{K}\) for every \(i \in I''\) ; assertion b) follows.

PROPOSITION 7. --- Let L be a commutative algebra of finite degree over K. The following properties are equivalent:

\begin{enumerate}
\def\labelenumi{(\roman{enumi})}
\item
  The algebra L is étale.
\item
  There exists a separable extension of K of finite degree that diagonalizes K.
\end{enumerate}

The implication \((\mathrm{ii})\Rightarrow(\mathrm{i})\) follows from V, §6, No.~3, p.~29, Proposition 2. Let us prove the implication \((\mathrm{i})\Rightarrow(\mathrm{ii})\) . Let \(\Omega\) be a separable closure of K. By Theorem 4 of V, §6, No.~7, p.~34, there exist extensions \(L_{1},\ldots,L_{n}\) of K of finite degree, contained in \(\Omega\) , such that L is isomorphic to the product

\(L_{1} \times \cdots \times L_{n}\) . Let N be a Galois extension of K that contains the \(L_{i}\) (V, §10, No.~1, p.~58), and let us prove that \(L_{(N)}\) is diagonalizable. By the primitive element theorem (V, §7, No.~4, p.~40, Theorem 1), for every \(i \in [1, n]\) , there exists an irreducible separable polynomial \(P_{i} \in K[X]\) such that \(L_{i}\) is isomorphic to \(K[X]/(P_{i})\) . Since N is a normal extension of K in which \(P_{i}\) admits a root, the polynomial \(P_{i}\) splits in N{[}X{]}, with simple roots. Consequently, the N-algebra \(L_{i(N)}\) , which is isomorphic to \(N[X]/(P_{i})\) , is isomorphic to \(N^{[L_{i}:K]}\) . Hence \(L_{(N)}\) is diagonalizable.

THEOREM 7. --- Let A be a central simple K-algebra of finite degree and L a subalgebra of A. The following properties are equivalent:

\begin{enumerate}
\def\labelenumi{(\roman{enumi})}
\item
  The algebra L is a maximal étale subalgebra of A.
\item
  There exist an extension \(K'\) of K, an integer \(n \geqslant 1\) , and an isomorphism \(\theta\) from \(\mathrm{A}_{(\mathrm{K}')}\) to \(\mathbf{M}_{n}(\mathrm{K}')\) that sends \(\mathrm{L}_{(\mathrm{K}')}\) to the set of diagonal matrices.
\item
  There exist \(K'\) , n, and \(\theta\) as in (ii), where the extension \(K'\) is, moreover, assumed Galois and of finite degree.
\end{enumerate}

It is clear that (iii) implies (ii).

If property (ii) holds, then \(L_{(K')}\) is a maximal commutative subalgebra of \(A_{(K')}\) (Proposition 6), and it is diagonalizable. The K-algebra L is then étale (V, §6, No.~3, p.~28, Definition 1), and it is a maximal commutative subalgebra of A (VIII, p.~261, Lemma 3, b)). We have proved that (ii) implies (i).

Suppose that property (i) holds. Since L is étale over K, by Proposition 7, there exists a Galois extension \(K_{1}\) of K of finite degree such that the \(K_{1}\) -algebra \(\mathrm{L}_{(\mathrm{K}_{1})}\) is diagonalizable. The algebra A is central simple; by (VIII, p.~252, Theorem 1), there exist a Galois extension \(K_{2}\) , a vector space V of finite dimension n over \(K_{2}\) , and an isomorphism \(\theta\) from \(\mathrm{A}_{(\mathrm{K}_{2})}\) to \(\mathrm{End}_{\mathrm{K}_{2}}(\mathrm{V})\) . By Proposition 1 of V, p.~55, we may assume \(K_{1}=K_{2}\) . By Proposition 4, b) of VIII, p.~264 and Lemma 3, b) of VIII, p.~261, \(\mathrm{L}_{(\mathrm{K}^{\prime})}\) is a maximal commutative subalgebra of \(\mathrm{A}_{(\mathrm{K}^{\prime})}\) , so \(\theta(\mathrm{L}_{(\mathrm{K}^{\prime})})\) is a maximal commutative subalgebra of \(\mathrm{End}_{\mathrm{K}^{\prime}}(\mathrm{V}^{\prime})\) . Let us apply Proposition 6 to the diagonalizable algebra \(\theta(\mathrm{L}_{(\mathrm{K}^{\prime})})\) : there exists a basis \((e_{1},\ldots,e_{n})\) of \(V^{\prime}\) over \(K^{\prime}\) such that \(\theta(\mathrm{L}_{(\mathrm{K}^{\prime})})\) consists of the endomorphisms of \(V^{\prime}\) with diagonal matrix with respect to this basis. So (i) implies (iii).

